% 2013 全国卷Ⅱ理科第 18 题：直三棱柱 ABC-A_1B_1C_1，D、E 为中点。
\begin{tikzpicture}[font=\footnotesize, line width=0.45pt]
  \coordinate (A) at (0,0);
  \coordinate (B) at (1.25,-1.25);
  \coordinate (C) at (4.25,0);
  \coordinate (A1) at (0,3.70);
  \coordinate (B1) at (2.15,2.00);
  \coordinate (C1) at (4.25,3.70);
  \coordinate (D) at ($(A)!0.5!(B)$);
  \coordinate (E) at ($(B)!0.5!(B1)$);
  % 可见棱和题中出现的截线。
  \draw (A)--(A1)--(C1)--(C)--cycle;
  \draw (A1)--(B1)--(C1) (B)--(B1);
  \draw (A1)--(D)--(C) (A1)--(E)--(C);
  % 透明棱与被遮挡的辅助关系。
  \draw[dashed] (A)--(B)--(C) (A)--(C) (A1)--(C) (D)--(E);
  \node[below left] at (A) {$A$};
  \node[below] at (B) {$B$};
  \node[below right] at (C) {$C$};
  \node[above left] at (A1) {$A_1$};
  \node[above] at (B1) {$B_1$};
  \node[above right] at (C1) {$C_1$};
  \node[below left] at (D) {$D$};
  \node[left] at (E) {$E$};
\end{tikzpicture}
